\documentclass[12pt,a4paper]{report}
\usepackage[utf8]{inputenc}
\usepackage[T1]{fontenc}
\usepackage[french]{babel}
\usepackage[top=2.5cm,bottom=2.5cm,left=2.8cm,right=2.8cm,headheight=15pt]{geometry}
\usepackage{xcolor,hyperref,titlesec,fancyhdr,booktabs,tabularx,enumitem,tcolorbox,lmodern,microtype,parskip,listings,float}
\tcbuselibrary{skins,breakable}

\definecolor{primary}{RGB}{30,58,95}
\definecolor{accent}{RGB}{37,99,235}
\definecolor{success}{RGB}{22,163,74}
\definecolor{warning}{RGB}{217,119,6}
\definecolor{danger}{RGB}{220,38,38}
\definecolor{codebg}{RGB}{30,30,30}
\definecolor{codefg}{RGB}{220,220,220}

\lstset{
  backgroundcolor=\color{codebg},basicstyle=\small\ttfamily\color{codefg},
  keywordstyle=\color{accent!80},commentstyle=\color{gray!60},
  stringstyle=\color{success!80},breaklines=true,frame=none,
  numbers=left,numberstyle=\tiny\color{gray!50},
  xleftmargin=12pt,aboveskip=6pt,belowskip=6pt
}

\newtcolorbox{infobox}[1][]{colback=blue!5,colframe=accent,fonttitle=\bfseries,title=#1,left=5pt,right=5pt,top=3pt,bottom=3pt,breakable}
\newtcolorbox{warnbox}[1][]{colback=orange!8,colframe=warning,fonttitle=\bfseries,title=#1,left=5pt,right=5pt,top=3pt,bottom=3pt,breakable}
\newtcolorbox{dangerbox}[1][]{colback=red!5,colframe=danger,fonttitle=\bfseries,title=#1,left=5pt,right=5pt,top=3pt,bottom=3pt,breakable}
\newtcolorbox{tipbox}[1][]{colback=green!5,colframe=success,fonttitle=\bfseries,title=#1,left=5pt,right=5pt,top=3pt,bottom=3pt,breakable}

\pagestyle{fancy}\fancyhf{}
\fancyhead[L]{\small\textcolor{primary}{\textbf{École Les Génies}}}
\fancyhead[R]{\small\textcolor{gray}{Cahier de Maintenance}}
\fancyfoot[C]{\small\thepage}
\renewcommand{\headrulewidth}{0.4pt}

\titleformat{\chapter}[display]{\normalfont\Huge\bfseries\color{primary}}{\chaptertitlename\ \thechapter}{20pt}{\Huge}
\titleformat{\section}{\normalfont\Large\bfseries\color{primary}}{\thesection}{1em}{}
\titleformat{\subsection}{\normalfont\large\bfseries\color{accent}}{\thesubsection}{1em}{}

\hypersetup{colorlinks=true,linkcolor=accent,urlcolor=accent,
  pdftitle={Cahier de Maintenance — Portail Numérique École Les Génies}}

\begin{document}

%% PAGE DE TITRE
\begin{titlepage}\centering\vspace*{2cm}
  {\color{primary}\rule{\linewidth}{3pt}}\vspace{0.5cm}
  {\Huge\bfseries\color{primary}Portail Numérique Scolaire\par}\vspace{0.3cm}
  {\LARGE\color{gray}École Les Génies\par}\vspace{0.8cm}
  {\color{primary}\rule{\linewidth}{1.5pt}}\vspace{1cm}
  {\huge\bfseries Cahier de Maintenance\par}\vspace{0.4cm}
  {\large Guide technique du responsable informatique\par}\vspace{2cm}
  \begin{tabular}{ll}\toprule
    \textbf{Document}   & Cahier de maintenance \\
    \textbf{Version}    & 1.0 \\
    \textbf{Date}       & Juillet 2026 \\
    \textbf{Public}     & Administrateur système / Développeur \\
    \textbf{Confidentialité} & Strictement interne \\
  \bottomrule\end{tabular}
  \vfill{\color{primary}\rule{\linewidth}{3pt}}
  {\small\color{gray}Document confidentiel — Usage interne uniquement}
\end{titlepage}

\tableofcontents\clearpage

%% ══════════════════════════════════════════════════
\chapter*{Introduction}\addcontentsline{toc}{chapter}{Introduction}

Ce document est le \textbf{cahier de maintenance technique} du Portail Numérique de l'École Les Génies. Il est destiné exclusivement au personnel informatique et aux développeurs chargés d'assurer le bon fonctionnement, la mise à jour et la sécurité de la plateforme.

\begin{dangerbox}[Attention : Avertissement]
Ce document contient des informations sensibles (architecture, accès, procédures d'urgence). Il ne doit jamais être partagé avec des utilisateurs finaux.
\end{dangerbox}

%% ══════════════════════════════════════════════════
\chapter{Architecture Technique}

\section{Vue d'ensemble}

L'application est composée de trois couches distinctes :

\begin{center}
\begin{tabular}{lllll}\toprule
\textbf{Couche} & \textbf{Technologie} & \textbf{Version} & \textbf{Port} & \textbf{Rôle} \\\midrule
Frontend principal & Vite / React & 6.x & 5173 & Interface Admin \\
Portail Enseignant & Next.js & 14+ & 3001 & Interface Enseignant \\
Portail Parent    & Next.js & 14+ & 3002 & Interface Parent \\
Backend API       & Laravel & 13.x & 8000 & Serveur REST API \\
Base de données   & MySQL & 8.x & 3306 & Persistance des données \\
\bottomrule\end{tabular}
\end{center}

\section{Arborescence du projet}

\begin{lstlisting}[language=bash]
projet-BD/
  backend-app/          # API Laravel (PHP 8.3+)
    app/Http/Controllers/Legacy/   # Controleurs metier
    database/migrations/           # Migrations BDD
    routes/api.php                 # Routes API
    .env                           # Configuration locale
  fronted/              # Applications frontend (Node.js)
    src/                           # Interface Admin (Vite/React)
    teacher-page-design/           # Portail Enseignant (Next.js)
    admin-interface-design (2)/    # Portail Parent (Next.js)
    package.json                   # Scripts de demarrage
  docs/                 # Documentation (ce fichier)
\end{lstlisting}

\section{Flux de communication}

\begin{infobox}[Flux général]
Navigateur → Frontend (Vite/Next.js) → API REST (Laravel, port 8000) → Base de données MySQL (\texttt{ecole2026\_local})
\end{infobox}

Toutes les requêtes API passent par \texttt{http://localhost:8000/api/legacy/}. Le préfixe \texttt{legacy} indique que l'API s'appuie sur la base de données existante de l'école.

%% ══════════════════════════════════════════════════
\chapter{Installation et Déploiement}

\section{Prérequis}

\begin{tabularx}{\textwidth}{lXl}\toprule
\textbf{Logiciel} & \textbf{Rôle} & \textbf{Version min.} \\\midrule
PHP & Serveur backend Laravel & 8.3+ \\
Composer & Gestionnaire de paquets PHP & 2.x \\
Node.js & Exécution du frontend & 18+ \\
npm & Gestionnaire de paquets JS & 9+ \\
MySQL & Base de données & 8.x \\
Git & Contrôle de version & 2.x \\
\bottomrule\end{tabularx}

\section{Installation initiale}

\subsection{Backend (Laravel)}
\begin{lstlisting}[language=bash]
cd projet-BD/backend-app

# 1. Installer les dependances PHP
composer install

# 2. Copier la configuration
cp .env.example .env

# 3. Configurer le fichier .env
nano .env
# Renseigner : DB_HOST, DB_PORT, DB_DATABASE, DB_USERNAME, DB_PASSWORD

# 4. Generer la cle applicative
php artisan key:generate

# 5. Lancer les migrations
php artisan migrate

# 6. Demarrer le serveur
php artisan serve --port=8000
\end{lstlisting}

\subsection{Frontend (React / Next.js)}
\begin{lstlisting}[language=bash]
cd projet-BD/fronted

# 1. Installer les dependances
npm install

# 2. Builder les portails Next.js (a faire une seule fois)
npm run build:next

# 3. Lancer tous les serveurs (mode production leger)
npm run dev
\end{lstlisting}

\begin{infobox}[Scripts disponibles]
\begin{itemize}[leftmargin=*,noitemsep]
  \item \texttt{npm run dev} — Mode production (CPU léger, recommandé au quotidien)
  \item \texttt{npm run dev:hot} — Mode développement avec hot-reload complet
  \item \texttt{npm run build:next} — Rebuild des portails Next.js après modification
  \item \texttt{npm run build} — Build de l'interface admin Vite
\end{itemize}
\end{infobox}

\section{Configuration du fichier \texttt{.env}}

\begin{lstlisting}[language=bash]
APP_NAME="Portail Les Genies"
APP_ENV=production        # local en developpement
APP_DEBUG=false           # false en production !
APP_URL=http://localhost

DB_CONNECTION=mysql
DB_HOST=127.0.0.1
DB_PORT=3306
DB_DATABASE=ecole2026_local
DB_USERNAME=root
DB_PASSWORD=votre_mot_de_passe

SESSION_DRIVER=database
LOG_LEVEL=error           # debug en developpement
\end{lstlisting}

\begin{dangerbox}[Sécurité]
Ne jamais mettre \texttt{APP\_DEBUG=true} en production. Cela expose la structure interne du code dans les messages d'erreur.
\end{dangerbox}

%% ══════════════════════════════════════════════════
\chapter{Base de Données}

\section{Connexion et informations générales}

\begin{tabularx}{\textwidth}{lX}\toprule
\textbf{Paramètre} & \textbf{Valeur} \\\midrule
Système & MySQL 8.x \\
Base de données & \texttt{ecole2026\_local} \\
Hôte & \texttt{127.0.0.1} (local) \\
Port & \texttt{3306} \\
Encodage & \texttt{utf8mb4} \\
\bottomrule\end{tabularx}

\section{Tables principales}

\begin{tabularx}{\textwidth}{lX}\toprule
\textbf{Table} & \textbf{Description} \\\midrule
\texttt{users} & Comptes utilisateurs (admin, enseignants, parents) \\
\texttt{people} & Personnes physiques (état civil) \\
\texttt{teachers} & Enseignants, lien vers \texttt{people} \\
\texttt{parents} & Parents d'élèves, lien vers \texttt{people} \\
\texttt{students} & Élèves, lien vers \texttt{people} \\
\texttt{subjects} & Matières enseignées \\
\texttt{school\_years} & Années académiques \\
\texttt{terms} & Trimestres \\
\texttt{rooms} & Salles de classe \\
\texttt{classes} & Classes (CP, CE1, etc.) \\
\texttt{enrollments} & Inscriptions des élèves dans les classes \\
\texttt{teacher\_class\_assignments} & Affectations titulaires aux classes \\
\texttt{grades} & Notes des élèves \\
\texttt{attendances} & Présences et absences \\
\texttt{disciplines} & Sanctions disciplinaires \\
\texttt{payments} & Paiements des frais de scolarité \\
\texttt{messages} & Messages entre admin et parents \\
\bottomrule\end{tabularx}

\section{Sauvegardes de la base de données}

\subsection{Sauvegarde manuelle}
\begin{lstlisting}[language=bash]
# Sauvegarde complete
mysqldump -u root -p ecole2026_local > backup_$(date +%Y%m%d_%H%M%S).sql

# Restauration
mysql -u root -p ecole2026_local < backup_YYYYMMDD_HHMMSS.sql
\end{lstlisting}

\subsection{Sauvegarde automatique (cron)}
\begin{lstlisting}[language=bash]
# Editer la crontab
crontab -e

# Ajouter la ligne suivante (sauvegarde quotidienne a 2h00)
0 2 * * * mysqldump -u root -pMOT_DE_PASSE ecole2026_local \
  > /backups/ecole_$(date +\%Y\%m\%d).sql
\end{lstlisting}

\begin{tipbox}[Bonne pratique]
Conserver les sauvegardes sur un disque externe ou un espace cloud. Tester périodiquement la restauration pour s'assurer de l'intégrité des sauvegardes.
\end{tipbox}

%% ══════════════════════════════════════════════════
\chapter{API Backend}

\section{Structure des routes}

Toutes les routes sont définies dans \texttt{backend-app/routes/api.php} sous le préfixe \texttt{/api/legacy/}.

\begin{tabularx}{\textwidth}{llX}\toprule
\textbf{Méthode} & \textbf{Endpoint} & \textbf{Description} \\\midrule
POST & \texttt{/auth/login-parent} & Connexion parent \\
POST & \texttt{/auth/login-teacher} & Connexion enseignant \\
POST & \texttt{/auth/login-admin} & Connexion administrateur \\
GET  & \texttt{/parent/dashboard} & Données tableau de bord parent \\
GET  & \texttt{/parent/notes/\{matricule\}} & Notes d'un élève \\
GET  & \texttt{/parent/bulletins/\{matricule\}} & Liste des bulletins \\
GET  & \texttt{/parent/bulletin-detail/\{matricule\}/\{idTrimes\}} & Détail bulletin \\
GET  & \texttt{/teacher/dashboard} & Données enseignant \\
GET  & \texttt{/admin/dashboard} & Statistiques admin \\
GET  & \texttt{/admin/eleves} & Liste des élèves \\
GET  & \texttt{/admin/paiements} & Liste des paiements \\
\bottomrule\end{tabularx}

\section{Contrôleurs (Controllers)}

Les contrôleurs sont dans \texttt{app/Http/Controllers/Legacy/} :

\begin{tabularx}{\textwidth}{lX}\toprule
\textbf{Fichier} & \textbf{Responsabilité} \\\midrule
\texttt{LegacyAuthController} & Authentification des 3 types d'utilisateurs \\
\texttt{LegacyParentDashboardController} & Données du portail parent (enfants, notes, bulletins) \\
\texttt{LegacyTeacherController} & Données du portail enseignant \\
\texttt{LegacyStudentController} & Gestion des élèves \\
\texttt{LegacyDashboardController} & Statistiques du tableau de bord admin \\
\texttt{LegacyPaymentController} & Gestion des paiements et impayés \\
\texttt{LegacyDisciplineController} & Gestion des sanctions \\
\texttt{LegacyMessageController} & Messagerie interne \\
\texttt{LegacyUserController} & Gestion des comptes utilisateurs \\
\texttt{LegacyStructureController} & Classes, salles, titulaires \\
\texttt{LegacyCoursController} & Matières et emploi du temps \\
\bottomrule\end{tabularx}

%% ══════════════════════════════════════════════════
\chapter{Gestion des Processus et Performances}

\section{Démarrage des serveurs}

\begin{lstlisting}[language=bash]
# Terminal 1 : Backend Laravel
cd projet-BD/backend-app
php artisan serve --port=8000

# Terminal 2 : Frontends (tous les 3 via concurrently)
cd projet-BD/fronted
npm run dev
\end{lstlisting}

\section{Problème de surcharge CPU}

\begin{warnbox}[Symptôme connu : surcharge CPU]
Si la machine est lente et que \texttt{htop} montre un \textit{Load Average} élevé ($>$ 4), c'est souvent dû à plusieurs serveurs \texttt{next dev} qui tournent simultanément.
\end{warnbox}

\textbf{Solution immédiate :}
\begin{lstlisting}[language=bash]
# Tuer tous les processus Next.js et Vite
pkill -f "next-server" ; pkill -f "next dev" ; pkill -f "vite"

# Relancer proprement
cd projet-BD/fronted && npm run dev
\end{lstlisting}

\textbf{Explication :} Le script \texttt{npm run dev} utilise désormais \texttt{next start} (mode production, $\sim$1\% CPU) au lieu de \texttt{next dev} (mode développement, $\sim$30--50\% CPU par serveur). Si des modifications sont apportées aux portails Next.js, il faut rebuilder avant de relancer :

\begin{lstlisting}[language=bash]
npm run build:next   # Rebuild les portails
npm run dev          # Relancer
\end{lstlisting}

\section{Vérification de l'état des services}

\begin{lstlisting}[language=bash]
# Verifier les processus actifs
ps aux | grep -E "artisan|next|vite" | grep -v grep

# Verifier les ports utilises
ss -tlnp | grep -E "5173|3001|3002|8000"

# Verifier la connexion a la base de donnees
mysql -u root -p -e "SHOW DATABASES;" | grep ecole
\end{lstlisting}

%% ══════════════════════════════════════════════════
\chapter{Sécurité}

\section{Authentification}

\begin{itemize}[leftmargin=*]
  \item Les mots de passe sont hachés avec \textbf{bcrypt} (12 rounds) via Laravel.
  \item L'espace administrateur est protégé par un \textbf{code d'accès préalable} avant la saisie des identifiants.
  \item Les sessions utilisateurs sont gérées côté backend via \texttt{localStorage} (token JWT ou identifiant de session).
\end{itemize}

\section{Gestion des rôles}

\begin{tabularx}{\textwidth}{lX}\toprule
\textbf{Rôle} & \textbf{Accès autorisés} \\\midrule
\texttt{directeur} / \texttt{administration} & Tout le module académique, élèves, pédagogie, suivi \\
\texttt{fondateur} & Idem directeur + gestion utilisateurs + configuration \\
\texttt{intendant} & Tableau de bord, paiements, impayés, configuration \\
\texttt{root} & Gestion utilisateurs, configuration système \\
\texttt{enseignant} & Portail enseignant uniquement (port 3001) \\
\texttt{parent} & Portail parent uniquement (port 3002) \\
\bottomrule\end{tabularx}

\begin{dangerbox}[Compte Root]
Le compte \texttt{root} a accès à la gestion des utilisateurs et à la configuration système. Son mot de passe doit être unique, complexe ($\geq$ 16 caractères) et connu d'une seule personne de confiance.
\end{dangerbox}

\section{Recommandations de sécurité}

\begin{enumerate}[leftmargin=*]
  \item \textbf{Sauvegardes chiffrées} : Chiffrer les dumps MySQL avec GPG avant stockage distant.
  \item \textbf{HTTPS} : En production, utiliser un reverse proxy (Nginx) avec certificat SSL.
  \item \textbf{Mises à jour} : Maintenir Laravel et Node.js à jour pour les correctifs de sécurité.
  \item \textbf{Logs} : Surveiller \texttt{backend-app/storage/logs/laravel.log} régulièrement.
  \item \textbf{Firewall} : En production, n'exposer que les ports nécessaires (80/443).
\end{enumerate}

%% ══════════════════════════════════════════════════
\chapter{Procédures de Maintenance Courante}

\section{Mise à jour des dépendances}

\subsection{Backend (PHP/Laravel)}
\begin{lstlisting}[language=bash]
cd backend-app
composer update                   # Met a jour les paquets PHP
php artisan migrate               # Applique les nouvelles migrations
php artisan config:cache          # Vide le cache de configuration
php artisan route:cache           # Vide le cache des routes
\end{lstlisting}

\subsection{Frontend (Node.js)}
\begin{lstlisting}[language=bash]
cd fronted
npm update                        # Met a jour les paquets JS
npm run build:next                # Rebuild si paquets Next.js modifies
\end{lstlisting}

\section{Réinitialisation du mot de passe d'un utilisateur}

Via l'interface (recommandé) : Menu \textbf{Utilisateurs} → sélectionner l'utilisateur → \textbf{Réinitialiser MDP}.

Via la base de données (en cas d'urgence) :
\begin{lstlisting}[language=bash]
# Generer un hash bcrypt (depuis la console Laravel)
cd backend-app
php artisan tinker
>>> Hash::make('NouveauMotDePasse123!')

# Appliquer dans la BDD
mysql -u root -p ecole2026_local -e \
  "UPDATE users SET password='HASH_GENERE' WHERE username='identifiant';"
\end{lstlisting}

\section{Ajout d'un nouvel enseignant ou parent}

\begin{enumerate}[leftmargin=*]
  \item Se connecter en tant qu'admin (Fondateur ou Root).
  \item Aller dans \textbf{Utilisateurs} → \textbf{Nouvel utilisateur}.
  \item Choisir le rôle (Enseignant ou Parent), remplir les informations.
  \item Communiquer le mot de passe temporaire généré à l'utilisateur.
  \item L'utilisateur se connecte et doit changer son mot de passe.
\end{enumerate}

\section{Gestion des logs Laravel}

\begin{lstlisting}[language=bash]
# Consulter les derniers logs
tail -100 backend-app/storage/logs/laravel.log

# Vider les logs (si trop volumineux)
> backend-app/storage/logs/laravel.log
\end{lstlisting}

%% ══════════════════════════════════════════════════
\chapter{Résolution de Problèmes Fréquents}

\begin{tabularx}{\textwidth}{lXl}\toprule
\textbf{Problème} & \textbf{Cause probable} & \textbf{Solution} \\\midrule
API répond 500 & Erreur PHP / table manquante & Vérifier \texttt{laravel.log} \\
\addlinespace
Page blanche sur le portail & Serveur frontend arrêté & Relancer \texttt{npm run dev} \\
\addlinespace
Connexion refusée (port 8000) & Serveur Laravel arrêté & Relancer \texttt{php artisan serve} \\
\addlinespace
Machine lente / CPU $>$ 80\% & Trop de serveurs \texttt{next dev} & \texttt{pkill -f next-server} puis \texttt{npm run dev} \\
\addlinespace
Erreur ``Table not found'' & Migration non appliquée & \texttt{php artisan migrate} \\
\addlinespace
Erreur CORS & URL backend incorrecte & Vérifier \texttt{VITE\_API\_URL} dans \texttt{.env} frontend \\
\addlinespace
Bulletin vide & Pas de notes saisies & Vérifier les notes dans le module Notes \& Moyennes \\
\addlinespace
Login impossible & MDP incorrect ou BDD hors ligne & Vérifier BDD + reset MDP si nécessaire \\
\bottomrule\end{tabularx}

\section{Logs et débogage}

\begin{lstlisting}[language=bash]
# Logs Laravel (backend)
tail -f backend-app/storage/logs/laravel.log

# Logs Node.js (frontend) - visibles dans le terminal de demarrage
npm run dev  # Les erreurs s'affichent directement

# Mode debug Laravel (developpement uniquement)
# Dans .env : APP_DEBUG=true, puis :
php artisan config:clear
\end{lstlisting}

%% ══════════════════════════════════════════════════
\chapter{Contacts et Responsabilités}

\begin{tabularx}{\textwidth}{llX}\toprule
\textbf{Rôle} & \textbf{Contact} & \textbf{Responsabilité} \\\midrule
Développeur principal & dev@lesgenies.cm & Maintenance code, déploiements \\
Administrateur BDD & dba@lesgenies.cm & Sauvegardes, intégrité données \\
Direction           & direction@lesgenies.cm & Validation des mises à jour majeures \\
\bottomrule\end{tabularx}

\vspace{1cm}
\begin{tipbox}[En cas d'urgence]
Si l'application est totalement inaccessible, la procédure est : (1) vérifier que MySQL tourne, (2) redémarrer les serveurs avec \texttt{pkill -f "next-server" \&\& pkill -f vite \&\& pkill -f artisan}, puis relancer le backend et le frontend.
\end{tipbox}

\vfill
\begin{center}
  \color{gray}\small\textit{École Les Génies — Portail Numérique v1.0 — Juillet 2026}\\
  \textit{Document de maintenance — Diffusion strictement restreinte}
\end{center}

\end{document}
